\newpage
\section{Introduction}

Avec l’augmentation constante du volume de documents mis en ligne, le plagiat est devenu une problématique centrale dans de nombreux domaines, notamment en milieu académique et juridique. Identifier l’origine d’un document ou d’une partie de texte est un défi important pour assurer l’intégrité des contenus et lutter contre la fraude intellectuelle.

L’objectif de ce projet est de mettre en œuvre des techniques de traitement automatique du langage naturel (\textit{NLP}) afin d’analyser le contenu textuel et d’évaluer la similarité entre documents. Cette approche repose sur plusieurs étapes clés : la transformation des textes en représentations exploitables, l’implémentation de métriques de comparaison, et l’utilisation de modèles d’apprentissage automatique pour classer et regrouper les textes selon leur origine.

Dans ce cadre, nous avons choisi de développer notre propre pipeline de détection, sans recourir à des implémentations toutes faites. Cela inclut la création d’un tokenizer, la vectorisation des documents, ainsi que l’implémentation et l’évaluation de différents algorithmes de classification. Une analyse quantitative des résultats sera menée afin de comparer les performances des différentes approches et d’identifier les facteurs influençant leur efficacité.

Ce rapport détaille les choix méthodologiques adoptés, les techniques mises en œuvre et les résultats obtenus, dans le but de mieux comprendre les enjeux liés à la détection de plagiat et d’évaluer l’efficacité des différentes stratégies d’analyse.
